% Generated draft: requires mathematical and source review.
\documentclass[11pt]{article}
\usepackage[margin=25mm]{geometry}
\usepackage{fontspec}
\setmainfont{Latin Modern Roman}
\usepackage{amsmath,amssymb,mathtools,unicode-math}
\setmathfont{Latin Modern Math}
\usepackage{enumitem,needspace,xurl,hyperref}
\setlength{\parindent}{0pt}
\setlength{\parskip}{5pt}
\setlength{\emergencystretch}{3em}
\newcommand{\N}{\mathbb{N}}
\newcommand{\Z}{\mathbb{Z}}
\newcommand{\Q}{\mathbb{Q}}
\newcommand{\R}{\mathbb{R}}
\newcommand{\C}{\mathbb{C}}
\newcommand{\F}{\mathbb{F}}
\newcommand{\A}{\mathbb{A}}
\newcommand{\PP}{\mathbb{P}}
\newcommand{\E}{\mathbb{E}}
\newcommand{\eps}{\varepsilon}
\newcommand{\dd}{\,\mathrm{d}}
\begin{document}
\section*{UnsolvedMath ID: 701567}
\textbf{Source problem code:} catalogue entry
\subsection{Definitions and mathematical statement}
Let \(\Omega\) be a nonempty finite set and let \(G\leq\operatorname{Sym}(\Omega)\) act transitively. Write
\[
\binom{\Omega}{2}=\{\{\alpha,\beta\}:\alpha,\beta\in\Omega,\ \alpha\ne\beta\}.
\]
The induced action on this set is \(g\{\alpha,\beta\}=\{g\alpha,g\beta\}\). For an orbit \(O\subseteq\binom{\Omega}{2}\), a proper block of imprimitivity is a nonempty subset \(B\subsetneq O\) such that, for every \(g\in G\), either \(gB=B\) or \(gB\cap B=\varnothing\). In particular, singleton blocks are included. Define the undirected graph
\[
\Gamma_{O,B}=(\Omega,O\setminus B).
\]
The group \(G\), in its specified action on \(\Omega\), has the \emph{road closure property} if \(\Gamma_{O,B}\) is connected for every such orbit \(O\) and every such proper block \(B\). Connectivity means that every two vertices can be joined by a finite path; paths of length zero are allowed. No lower bound on \(|\Omega|\) is supplied. If there are no applicable orbits or proper blocks, the defining universal condition is vacuous.

A transitive action is primitive if its only invariant partitions are the partition into singletons and the one-part partition. Following the source's description, call the action \emph{basic} if it is not contained, up to identification of the underlying sets, in a wreath product in a nontrivial product action. Explicitly, this excludes an identification \(\Omega\cong\Delta^k\), with \(|\Delta|\geq2\) and \(k\geq2\), under which \(G\) is contained in \(\operatorname{Sym}(\Delta)\wr\operatorname{Sym}(k)\) acting by coordinate permutations and independent permutations within coordinates.

The classification problem is to determine all finite basic primitive permutation groups \(G\leq\operatorname{Sym}(\Omega)\) for which there exist an orbit \(O\) on \(\binom{\Omega}{2}\) and a nonempty proper block \(B\) for the action on \(O\) such that
\[
(\Omega,O\setminus B)\quad\text{is disconnected}.
\]

The source further asks whether every group in this classification either has a subgroup of index at most \(3\) or is almost simple. Here almost simple means that, for some finite nonabelian simple group \(S\), one has \(S\leq G\leq\operatorname{Aut}(S)\), up to group isomorphism. The subgroup clause has an unresolved qualification: as literally written, \(\exists H\leq G\) with \([G:H]\leq3\) allows \(H=G\), so it is automatic. If a proper subgroup is intended, the substantive question is whether every group in the classification satisfies
\[
\bigl(\exists H<G:\ [G:H]\in\{2,3\}\bigr)
\quad\text{or}\quad
\bigl(G\text{ is almost simple}\bigr).
\]
Properness is not stated in the supplied record, so this latter formulation is recorded as a possible intended reading, not an unqualified replacement.

As background, the source states that a group with the road closure property must be primitive and basic and cannot have an imprimitive subgroup of index \(2\). It also excludes permutation groups arising from triality whose socle is denoted \(D_4(q)\) and whose intersection with a point stabiliser is the parabolic subgroup corresponding to the three leaves of the Coxeter--Dynkin diagram for \(D_4\). These are supplied background assertions, not results established here. The record does not specify the parameter restrictions or a complete construction of this triality family; the parabolic description is retained without choosing an additional indexing convention.
\subsection{Short English statement}
Classify the finite basic primitive permutation groups for which deleting a proper block from some orbit on unordered pairs leaves a disconnected graph on the original point set. The source also asks whether every such group has a subgroup of index at most \(3\) or is almost simple; whether the subgroup must be proper is unresolved. The supplied registry labels the problem open, without independent verification here.
\subsection{Sources}
\begin{itemize}
\item Supplied local UnsolvedMath registry record, permanent ID \(701567\), exact source code \texttt{catalogue entry}, title \emph{Kourovka Notebook Problem 19.19}. Upstream ID: \(701567\); upstream revision: \texttt{d371ee917ac9647fbb06cd6fc529deec6f3dda82}. The record supplies the status \texttt{open} and the update timestamp \texttt{2026-10-16T14:00:00Z}; these are registry metadata, not independently verified mathematical or literature assessments.
\item P. J. Cameron, proposer, \(2018\). Attributed source: \emph{The Kourovka Notebook (Unsolved Problems in Group Theory)}, original problem \(19.19\), nineteenth issue (\(2018\)); the record additionally identifies the twenty-first issue edition (\(2026\)), \texttt{arXiv:1401.0300v48}.
\item The supplied collection description names E. I. Khukhro and V. D. Mazurov as editors of the Kourovka Notebook collection covering Issues \(1\)--\(21\) (\(1965\)--\(2026\)). No source was independently visited or checked for this draft.
\end{itemize}
\end{document}
